\section{Propositional resolution}
\label{sec:prop-metatheory-04}

The clause object in this section is exactly the finite set defined in
\cref{sec:prop-syntax-semantics-04}.  A valuation satisfies a clause set \(S\) iff it
satisfies every \(C\in S\), so the set is interpreted conjunctively.  The
complement of \(p\) is \(\neg p\), and the complement of \(\neg p\) is
\(p\).  A clause containing both a literal and its complement is
\emph{tautological}; its formula interpretation is valid, so deleting it from
a clause set preserves the satisfying valuations of that set.

\begin{definition}[Resolvent]
Let \(E\) and \(F\) be clauses, let \(L\in E\), and let its complement
\(\overline L\) belong to \(F\).  Their \emph{resolvent on \(L\)} is the
clause
\[
  (E\setminus\set{L})\cup(F\setminus\set{\overline L}).
\]
A \emph{resolution derivation} from a clause set \(S\) is a finite sequence
whose members are clauses in \(S\) or resolvents of earlier members.
\end{definition}

\begin{theorem}[Soundness of propositional resolution]
Every valuation satisfying both parent clauses satisfies their resolvent.
Consequently, deriving the empty clause proves that the initial clause set is
not satisfiable.
\end{theorem}
\begin{proof}
Let \(E,F\) be the parents and let
\(R=(E\setminus\set{L})\cup(F\setminus\set{\overline L})\).  Suppose a
valuation falsifies \(R\).  It therefore falsifies every nonpivot literal in
both parents.  To satisfy \(E\), it must satisfy \(L\); to satisfy \(F\), it
must satisfy \(\overline L\).  No valuation satisfies complementary
literals.  Therefore no valuation can satisfy both parents while falsifying
the resolvent.  Induction on a resolution derivation preserves satisfaction.
The empty clause is false under every valuation.
\end{proof}

\begin{theorem}[Refutational completeness of propositional resolution]
\label{thm:prop-resolution-complete}
Let \(S\) be a finite set of clauses.  Then \(S\) is not satisfiable
iff the empty clause has a resolution derivation from \(S\).
\end{theorem}
\begin{proof}
Soundness proves the reverse implication.  For the forward implication,
delete tautological clauses.  Induct on the number of variables occurring in
\(S\).  If there are no variables, an unsatisfiable clause set contains the
empty clause.

Choose a variable \(p\).  Form \(S_p\) from all clauses not containing
\(p\) or \(\neg p\), together with every non-tautological resolvent on
\(p\) of a clause containing \(p\) and a clause containing \(\neg p\).  We claim that
\(S\) is satisfiable iff \(S_p\) is satisfiable.  Restricting a valuation
satisfying \(S\) to the other variables satisfies all retained clauses and all
resolvents by soundness.  Conversely, let \(w\) satisfy \(S_p\).  A clause
\(\set{p}\cup A\), where \(p\notin A\), whose remainder \(A\) is false
under \(w\) requires \(p=\True\); a clause
\(\set{\neg p}\cup B\), where \(\neg p\notin B\), whose remainder \(B\)
is false requires \(p=\False\).  Both requirements cannot occur:
otherwise the resolvent \(A\cup B\) on \(p\) would be false under
\(w\); being false under \(w\), it contains no complementary pair of
literals, so it is not tautological, belongs to \(S_p\), and contradicts
the assumption that \(w\) satisfies \(S_p\).  Choose either
permitted value for \(p\).  The resulting extension satisfies \(S\), proving
the claim.

If \(S\) is not satisfiable, neither is \(S_p\).  It contains fewer variables,
so the induction hypothesis gives a resolution derivation of the empty clause
from \(S_p\).  Every noninitial clause of \(S_p\) is itself a resolvent of
clauses in \(S\).  Splicing those one-step derivations into the derivation from
\(S_p\) gives a resolution derivation from \(S\).
\end{proof}

\begin{remark}[Attribution and scope of the finiteness hypothesis]
\label{rem:resolution-finiteness}
The variable-elimination step in the proof of
\cref{thm:prop-resolution-complete}---passing from \(S\) to \(S_p\) and
discarding the variable \(p\)---is the elimination rule of Davis and
Putnam \cite{DavisPutnam1960}.  The finiteness hypothesis on \(S\) is
consumed by the induction, not by the result: an infinite clause set may
contain infinitely many variables, and the induction on the number of
variables then has no measure to decrease.  The restriction is removed by
propositional compactness.  Suppose an arbitrary clause set \(S\) is not
satisfiable.  A valuation satisfies a clause \(C\) iff it satisfies the
formula \(\operatorname{form}_{\CPL}(C)\) of
\cref{sec:prop-syntax-semantics-04}, so the formula set
\(\set{\operatorname{form}_{\CPL}(C):C\in S}\) is not satisfiable, and
\cref{thm:prop-compactness} gives a finite subset of it that is not
satisfiable.  Selecting for each of its finitely many members one clause
of \(S\) interpreting it yields a finite \(S_0\subseteq S\) that is not
satisfiable.  \Cref{thm:prop-resolution-complete} applied to \(S_0\)
gives a resolution derivation of the empty clause from \(S_0\), and every
resolution derivation from \(S_0\) is one from \(S\).  Since the
soundness theorem above is stated without finiteness of the clause set,
the equivalence of \cref{thm:prop-resolution-complete} therefore holds
for every clause set, finite or not.
\end{remark}

Resolution is a refutation calculus: it derives the empty clause from an
unsatisfiable clause set.  Conversion to conjunctive normal form and
resolution must therefore be kept distinct.  The former supplies clauses; the
latter derives clauses.  The original resolution method is due to Robinson
\cite{Robinson1965}; normal-form transformations and their preservation
properties are treated systematically in \cite{BaazEglyLeitsch2001}.

